\documentclass[border=0pt]{standalone}
\usepackage[T1]{fontenc}
\usepackage[utf8]{inputenc}
\usepackage[spanish,es-noshorthands,es-nodecimaldot]{babel}
\usepackage{amsmath,amssymb}
\usepackage{newpxtext,newpxmath}
\usepackage{xcolor}
\usepackage{tikz}
\usetikzlibrary{arrows.meta,positioning,fit,backgrounds,calc}
\definecolor{navy}{HTML}{17243A}
\definecolor{blue}{HTML}{2A5CAA}
\definecolor{orange}{HTML}{D2691E}
\definecolor{cream}{HTML}{FCEFE3}
\definecolor{pale}{HTML}{EEF2F8}
\definecolor{text}{HTML}{1E232B}
\definecolor{muted}{HTML}{5F6873}
\definecolor{rule}{HTML}{C8D0DA}
\definecolor{sky}{HTML}{9DB9E3}
\definecolor{redact}{HTML}{3C4552}

\newcommand{\no}[1]{{\color{orange}#1}}
\setlength{\parindent}{0pt}
\setlength{\parskip}{0.45em}
\begin{document}
\hyphenpenalty=10000\exhyphenpenalty=10000
\begin{minipage}{26.783cm}\fontsize{27}{33.5}\selectfont\raggedright 
Un hogar empieza en $x=(j,a,w)$; al año siguiente vemos qué decidió ($e$) y con qué
coche quedó ($\delta'$). Su contribución a la verosimilitud suma los caminos que no
vemos:
\[ \begin{aligned} \ell(\theta)=\;&s(x)\,\Pi\big(\delta'\mid j,\bar a\big)\,\mathbf{1}[D{=}1]\\
   &+\big(1-s(x)\big)\sum_{w'}\no{Z(w'\mid x)}\;\Pi\big(e,\delta'\mid j,a{+}1,w'\big). \end{aligned} \]
La reparación no observada entra como una \textbf{mezcla} de las dos campanas:
\[ \no{Z(w'\mid x)}=\sum_{r\in\{0,1\}}\Pi\big(r\mid x;\,P_m,R_m\big)\;
   f\big(w'\mid x,r\big). \]
Los pesos $\Pi(r\mid x)$ cambian con el precio de reparación $R_m$ de cada mercado;
las densidades $f$, no.

\textbf{Estimación por punto fijo anidado.} Para cada $\theta$ se resuelven la
Bellman de cada tipo de hogar y los precios de equilibrio $P_m(\theta)$ de cada
mercado, y se evalúa $LL(\theta)=\sum\log\ell(\theta)$, como en Rust (1987) y
Gillingham et al.\ (2022).
\end{minipage}
\end{document}
